% !TEX root = ../main.tex
\documentclass[../main.tex]{subfiles}

\begin{document}

\selectlanguage{english}

\section{CHAPTER 5: MACHINE LEARNING}

\subsection{5.1. Definition and concepts of machine learning}
Machine learning designs algorithms that allow software systems to optimize task execution performance autonomously through training data. Structural classification isolates Supervised Learning (mapping inputs directly to labeled fields), Unsupervised Learning (discovering latent patterns in unlabeled data), and Reinforcement Learning (optimizing dynamic behaviors via rewards and penalties).

\subsection{5.2. Decision tree learning}
The ID3 algorithm compiles descriptive categorization models by leveraging information-theoretic equations as a heuristic mechanism to select the absolute highest-utility splitting variable at each branch node.
\begin{itemize}
    \item \textbf{Information Entropy Expression (Data Impurity Measurement):} For dataset $S$ spanning $c$ target classes where $p_i$ is the density of class $i$:
    \begin{equation}
        Entropy(S) = - \sum_{i=1}^{c} p_i \log_2(p_i)
    \end{equation}
    \item \textbf{Information Gain Expression (Attribute $A$ Splitting Value):} Computes the expected reduction in entropy achieved by partitioning dataset $S$ along attribute $A$:
    \begin{equation}
        Gain(S, A) = Entropy(S) - \sum_{v \in Values(A)} \frac{|S_v|}{|S|} Entropy(S_v)
    \end{equation}
    Where $S_v$ isolates the subset of $S$ where attribute $A$ holds value $v$.
    \item \textbf{Core Heuristic Selector Equation:} Branch node creation is governed by maximizing information extraction:
    \begin{equation}
        A^* = \arg\max_{A} Gain(S, A)
    \end{equation}
\end{itemize}

\subsection{5.3. Naive Bayes classification}
The Naive Bayes engine constructs classification decisions by enforcing the **Naive Independence Assumption**, asserting that all descriptive attributes $X = (x_1, x_2, ..., x_d)$ are conditionally independent given the targeted class variable $C_k$.

\textbf{Algorithmic Characterization Formula:} Imposing independence factors multi-variable distributions into basic atomic chain products:
\begin{equation}
    P(C_k | x_1, ..., x_d) = \frac{P(C_k) \cdot P(x_1, ..., x_d | C_k)}{P(x_1, ..., x_d)} \propto P(C_k) \cdot \prod_{j=1}^{d} P(x_j | C_k)
\end{equation}
The final classification is isolated by selecting the class that maximizes the structural Maximum A Posteriori (MAP) criteria:
\begin{equation}
    \hat{y} = \arg\max_{C_k} \left[ P(C_k) \cdot \prod_{j=1}^{d} P(x_j | C_k) \right]
\end{equation}
The mathematical chain product $\prod P(x_j | C_k)$ allows the model to efficiently compute boundaries over sparse datasets, entirely bypassing the curse of dimensionality.

\subsection{5.4. K-nearest neighbor}
KNN is a non-parametric, instance-based lazy learning framework. The model skips parametric training entirely, meaning its decision boundaries are defined explicitly by spatial geometry calculations over vector fields.
\begin{itemize}
    \item \textbf{Generalized Minkowski Metric Expression ($L_p$ norm):} Quantifies spatial proximity configurations between a query sample $x$ and training record $y$ across a $d$-dimensional space:
    \begin{equation}
        D_p(x, y) = \left( \sum_{j=1}^{d} |x_j - y_j|^p \right)^{\frac{1}{p}}
    \end{equation}
    Modulating parameter $p$ modifies geometric features: setting $p=2$ yields linear *Euclidean* metrics ($L_2$ norm); setting $p=1$ maps grid-like *Manhattan* distance pathways ($L_1$ norm).
    \item \textbf{Distance-Weighted Label Assignment Selector:} Let $N_K(x)$ represent the isolated set of $K$ nearest neighbors computed via metric $D_p$. To optimize classification boundaries, closer records are assigned higher voting weights via an inverse relation:
    \begin{equation}
        \hat{y} = \arg\max_{v} \sum_{y_i \in N_K(x)} w_i \cdot I(v_i = v) \quad \text{where} \quad w_i = \frac{1}{D_p(x, y_i)^2 + \epsilon}
    \end{equation}
    Where $I(\cdot)$ is the standard binary indicator function and $\epsilon \rightarrow 0$ is a stabilizing constant preventing division by zero. This mathematical formulation yields highly adaptive, non-linear local decision fields (\textit{Voronoi tessellation}).
\end{itemize}

\subsection{5.5. Other machine learning algorithms}
The textbook by Russell \& Norvig extends machine learning into structures modeling biological neural networks through Artificial Neural Networks (ANNs).
\begin{itemize}
    \item \textbf{Mathematical Node Profiling (Perceptron):} An input vector $x$ undergoes a linear combination via a dot product with weight matrix $w$, adding a bias offset $b$, before processing through a non-linear activation operator ($g$ e.g., Sigmoid, Tanh, or ReLU):
    \begin{equation}
        h_w(x) = g\left( \sum_{i=1}^{d} w_i x_i + b \right) = g(w^T x + b)
    \end{equation}
    \item \textbf{The Backpropagation Optimization Solver:} The core optimization engine driving deep architectures. The algorithm leverages the calculus **Chain Rule** to compute error gradients, distributing weight corrections $\Delta w_{ji}$ backward from the output layer to hidden layers to minimize global structural loss values ($E$):
    \begin{equation}
        \Delta w_{ji} = -\alpha \frac{\partial E}{\partial w_{ji}} = -\alpha \frac{\partial E}{\partial a_j} \cdot \frac{\partial a_j}{\partial z_j} \cdot \frac{\partial z_j}{\partial w_{ji}}
    \end{equation}
    Where $\alpha$ dictates the network learning rate, $z_j$ defines the net linear activation at node $j$, and $a_j = g(z_j)$ marks the final activated output profile of that node.
\end{itemize}

\end{document}
